% Options for packages loaded elsewhere
\PassOptionsToPackage{unicode}{hyperref}
\PassOptionsToPackage{hyphens}{url}
\PassOptionsToPackage{dvipsnames,svgnames,x11names}{xcolor}
\documentclass[
  10pt,
  fontset=fandol]{ctexart}
\usepackage{xcolor}
\usepackage[margin=16mm]{geometry}
\usepackage{amsmath,amssymb}
\setcounter{secnumdepth}{-\maxdimen} % remove section numbering
\usepackage{iftex}
\ifPDFTeX
  \usepackage[T1]{fontenc}
  \usepackage[utf8]{inputenc}
  \usepackage{textcomp} % provide euro and other symbols
\else % if luatex or xetex
  \usepackage{unicode-math} % this also loads fontspec
  \defaultfontfeatures{Scale=MatchLowercase}
  \defaultfontfeatures[\rmfamily]{Ligatures=TeX,Scale=1}
\fi
\usepackage{lmodern}
\ifPDFTeX\else
  % xetex/luatex font selection
\fi
% Use upquote if available, for straight quotes in verbatim environments
\IfFileExists{upquote.sty}{\usepackage{upquote}}{}
\IfFileExists{microtype.sty}{% use microtype if available
  \usepackage[]{microtype}
  \UseMicrotypeSet[protrusion]{basicmath} % disable protrusion for tt fonts
}{}
\makeatletter
\@ifundefined{KOMAClassName}{% if non-KOMA class
  \IfFileExists{parskip.sty}{%
    \usepackage{parskip}
  }{% else
    \setlength{\parindent}{0pt}
    \setlength{\parskip}{6pt plus 2pt minus 1pt}}
}{% if KOMA class
  \KOMAoptions{parskip=half}}
\makeatother
\setlength{\emergencystretch}{3em} % prevent overfull lines
\providecommand{\tightlist}{%
  \setlength{\itemsep}{0pt}\setlength{\parskip}{0pt}}
\usepackage{amsmath,amssymb,mathtools,needspace}
\xeCJKDeclareCharClass{CJK}{"2032,"0400->"04FF,"2160->"216B,"2460->"2473,"25A0,"25C7,"25CB,"25CF}
\IfFontExistsTF{Arial Unicode MS}{\xeCJKsetup{AutoFallBack=true}\setCJKfallbackfamilyfont{\CJKrmdefault}{Arial Unicode MS}}{}
\setcounter{secnumdepth}{0}
\setlength{\emergencystretch}{2em}
\allowdisplaybreaks
\usepackage{bookmark}
\IfFileExists{xurl.sty}{\usepackage{xurl}}{} % add URL line breaks if available
\urlstyle{same}
\hypersetup{
  pdftitle={追赶法},
  colorlinks=true,
  linkcolor={Maroon},
  filecolor={Maroon},
  citecolor={Blue},
  urlcolor={Blue},
  pdfcreator={LaTeX via pandoc}}

\title{追赶法}
\author{}
\date{}

\begin{document}
\maketitle

\subsubsection{7.4.3 追赶法}\label{ux8ffdux8d76ux6cd5}

在一些实际问题中，例如解常微分方程边值问题，解热传导方程以及船体数学放样中建立三次样条函数等中，都会要求解系数矩阵为对角占优的三对角方程组

\[
\begin{pmatrix}
b_1&c_1&&&&&\\
a_2&b_2&c_2&&&&\\
&\ddots&\ddots&\ddots&&&\\
&&a_i&b_i&c_i&&\\
&&&\ddots&\ddots&\ddots&\\
&&&&a_{n-1}&b_{n-1}&c_{n-1}\\
&&&&&a_n&b_n
\end{pmatrix}
\begin{pmatrix}x_1\\x_2\\\vdots\\x_i\\\vdots\\x_{n-1}\\x_n\end{pmatrix}
=
\begin{pmatrix}f_1\\f_2\\\vdots\\f_i\\\vdots\\f_{n-1}\\f_n\end{pmatrix},
\]

（方程组及说明续下页。）

\begin{quote}
校注（agent 补充，原书疑误）：式（7.4.11）步 4 的求和项原图印作 \(l_{ik}y_h\)，这里保留；求和指标为 \(k\) 且来源为 \(Ly=b\)，相应项应为 \(l_{ik}y_k\)。\href{../assets/ch07b/197-solves.png}{原式局部图}。
\end{quote}

\href{../../source-images/pdf-198.jpeg}{原页图像}

（续 7.4.3 追赶法。）

简记作

\[
Ax=f,
\tag{7.4.12}
\]

其中 \(A\) 满足下列对角占优条件：

（1）\(|b_1|>|c_1|>0\)；

（2）\(|b_i|\geqslant|a_i|+|c_i|\)，\(a_i,c_i\neq0\)（\(i=2,3,\cdots,n-1\)）；

（3）\(|b_n|>|a_n|>0\)。

下面利用矩阵的直接三角分解法来推导解三对角方程组（7.4.12）的计算公式。由系数阵 \(A\) 的特点，可以将 \(A\) 分解为两个三角阵的乘积，即

\[
A=LU
\]

其中 \(L\) 为下三角阵，\(U\) 为单位上三角阵。下面说明这种分解是可能的。设

\[
A=\begin{pmatrix}
b_1&c_1&&&\\
a_2&b_2&c_2&&\\
&\ddots&\ddots&\ddots&\\
&&a_{n-1}&b_{n-1}&c_{n-1}\\
&&&a_n&b_n
\end{pmatrix}
=\begin{pmatrix}
\alpha_1&&&&\\
\gamma_2&\alpha_2&&&\\
&\gamma_3&\alpha_3&&\\
&&\ddots&\ddots&\\
&&&\gamma_n&\alpha_n
\end{pmatrix}
\begin{pmatrix}
1&\beta_1&&&\\
&1&\beta_2&&\\
&&\ddots&\ddots&\\
&&&1&\beta_{n-1}\\
&&&&1
\end{pmatrix},
\tag{7.4.13}
\]

其中 \(\alpha_i,\beta_i,\gamma_i\) 为待定系数。比较式（7.4.13）两端即得

\[
\left.\begin{aligned}
b_1&=\alpha_1,\quad c_1=\alpha_1\beta_1,\\
a_i&=\gamma_i,\quad b_i=\gamma_i\beta_{i-1}+\alpha_i\quad(i=2,\cdots,n),\\
c_i&=\alpha_i\beta_i\quad(i=2,3,\cdots,n-1).
\end{aligned}\right\}
\tag{7.4.14}
\]

由 \(\alpha_1=b_1\neq0\)，\(|b_1|>|c_1|>0\)，\(\beta_1=c_1/b_1\)，得 \(0<|\beta_1|<1\)。下面用归纳法证明

\[
|\alpha_i|>|c_i|\neq0\quad(i=1,2,\cdots,n-1),
\tag{7.4.15}
\]

即

\[
0<|\beta_i|<1,
\]

从而由式（7.4.14）可求出 \(\beta_i\)。

式（7.4.15）对 \(i=1\) 是成立的。现设式（7.4.15）对 \(i-1\) 成立，求证对 \(i\) 亦成立。

由归纳法假设 \(0<|\beta_{i-1}|<1\)，又由式（7.4.15）及 \(A\) 的假设条件，有

\[
|\alpha_i|=|b_i-a_i\beta_{i-1}|\geqslant|b_i|-|a_i\beta_{i-1}|>|b_i|-|a_i|\geqslant|c_i|\neq0,
\]

也就是 \(0<|\beta_i|<1\)。由式（7.4.14）得到

\[
\alpha_i=b_i-a_i\beta_{i-1}\quad(i=2,3,\cdots,n),\quad
\beta_i=c_i/(b_i-a_i\beta_{i-1})\quad(i=2,3,\cdots,n-1).
\]

这就是说，由 \(A\) 的假设条件完全确定了 \(\{\alpha_i\},\{\beta_i\},\{\gamma_i\}\)，实现了 \(A\) 的 LU 分解。

求解 \(Ax=f\) 等价于解两个三角方程组 \(Ly=f\) 与 \(Ux=y\)，先后求 \(y\) 与 \(x\)，从而得到以下解三对角方程组的\textbf{追赶法公式}：

\textbf{步 1} 计算 \(\{\beta_i\}\) 的递推公式

\[
\beta_1=c_1/b_1,\quad\beta_i=c_i/(b_i-a_i\beta_{i-1})\quad(i=2,3,\cdots,n-1);
\]

\textbf{步 2} 解 \(Ly=f\)：

\[
y_1=f_1/b_1,\quad y_i=(f_i-a_iy_{i-1})/(b_i-a_i\beta_{i-1})\quad(i=2,3,\cdots,n);
\]

\textbf{步 3} 解 \(Ux=y\)：

\[
x_n=y_n,\quad x_i=y_i-\beta_ix_{i+1}\quad(i=n-1,n-2,\cdots,2,1).
\]

\href{../../source-images/pdf-199.jpeg}{原页图像}

（续 7.4.3 追赶法。）

将计算系数 \(\beta_1\to\beta_2\to\cdots\to\beta_{n-1}\) 及 \(y_1\to y_2\to\cdots\to y_n\) 的过程称为\textbf{追的过程}，将计算方程组的解 \(x_n\to x_{n-1}\to\cdots\to x_1\) 的过程称为\textbf{赶的过程}。

总结上述讨论，有以下定理。

\textbf{定理 7.9} 设有三对角方程组 \(Ax=f\)，其中 \(A\) 满足对角占优条件，则 \(A\) 为非奇异矩阵且由追赶法计算公式中 \(\{\alpha_i\},\{\beta_i\}\) 满足：

\(1^\circ\) \(0<|\beta_i|<1\)（\(i=1,2,\cdots,n-1\)）；

\(2^\circ\)

\[
0<|c_i|\leqslant|b_i|-|a_i|<|\alpha_i|<|b_i|+|a_i|\quad(i=2,3,\cdots,n-1),
\]

\[
0<|b_n|-|a_n|<|\alpha_n|<|b_n|+|a_n|.
\]

追赶法公式实际上就是把 Gauss 消去法用到求解三对角方程组上去的结果。这时由于 \(A\) 特别简单，因此使得求解的计算公式也非常简单，计算量仅为 \(5n-4\) 次乘除法运算，而另外增加解一个方程组 \(Ax=f_2\) 仅需增加 \(3n-2\) 次乘除运算。易见追赶法的计算量是比较小的。

定理 7.9 的 \(1^\circ\)、\(2^\circ\) 说明追赶法计算公式中不会出现中间结果数量级的巨大增长和舍入误差的严重累积。

在用计算机计算时，只需用三个一维数组分别存储 \(A\) 的三条对角线元素 \(\{a_i\},\{b_i\},\{c_i\}\)，此外还需要用两组存储单元保存 \(\{\beta_i\},\{y_i\}\) 或 \(\{x_i\}\)。

\subsection{本节覆盖原页的校注记录}\label{ux672cux8282ux8986ux76d6ux539fux9875ux7684ux6821ux6ce8ux8bb0ux5f55}

下列记录按原页关联，含同页相邻小节的校注；计算时同时读取上文校注和完整原页。

\begin{itemize}
\tightlist
\item
  \texttt{ch07b-E003}：原书 PDF 页 197
\end{itemize}

\end{document}
